\documentclass[10pt]{article}
\usepackage{template}
\usepackage{amsmath,amssymb,microtype,tikz,graphicx,xurl}
\usetikzlibrary{arrows.meta,positioning,calc}
\setlength{\emergencystretch}{3em}
\raggedbottom
\SetKeys{date=2026-10-30,vol=017,title={每日有趣新闻｜时政、科学与数学},author={ChatGPT 编辑部},note={\small 采编窗口：北京时间2026年10月28—30日；截稿：10月30日上午。时政区分官方声明与独立报道；数学预印本不等同于同行评议。}}
\renewcommand{\newsitem}[3]{\par\medskip{\large\bfseries\color{newsblue}#1\par}{\footnotesize\color{newsgray}#2\par}\smallskip\color{black}#3\par\medskip{\color{newsgray}\rule{\linewidth}{0.25pt}}\par}
\newcommand{\newslead}[3]{\par\medskip{\LARGE\bfseries\color{newsblue}#1\par}\vspace{1mm}{\small\color{newsgray}#2\par}\smallskip\color{black}#3\par\medskip{\color{newsgray}\rule{\linewidth}{0.45pt}}\par}
\newcommand{\newssub}[1]{\par\smallskip\noindent{\bfseries\color{newsblue}#1}\quad}
\begin{document}
\begin{documents}
\raggedcolumns
\newslead{二十届五中全会闭幕：军方人事调整与「全面从严治党」为何成为中心议题}{国内政治深度｜10月29日新华社授权公报；对照中央社、路透社与美联社}{

中国共产党第二十届中央委员会第五次全体会议于10月26日至29日在北京举行。29日公布的正式公报显示，会议审议通过《中共中央关于持之以恒推进全面从严治党若干重大问题的决定》，并增选中央军委副主席张升民为中央政治局委员。公报同时确认对多名高级党政军干部的党籍处分决定。与日常新闻发布不同，全会是中央委员会在党代会之间进行集体议事和作出重大政治决定的重要场合；理解它，需要把\textbf{正式作出的组织决定、政治口号和外部观察者的制度推断}分开。\cite{XinhuaPlenum,CNAPlenum,ReutersPlenum}

\newssub{先看确定发生了什么。}新华社公报记载，出席全会的中央委员162人、候补中央委员117人。会议听取和讨论了习近平受中央政治局委托所作的工作报告，审议通过关于持续推进全面从严治党的决定，并将张升民增选为政治局委员。张升民此前已任中央军委副主席，这次新增的是党内政治局职务；\textbf{中央军委职务与中央政治局席位并不是同一职位}，不能将这次人事变化写成他刚刚成为军委副主席。\cite{XinhuaPlenum,CNAPlenum}

公报还确认了涉及张又侠、刘振立、刘建超、马兴瑞等人的既有党籍处分。这里的「确认」有程序意义：它说明有关组织处理在全会层面获得确认，却不能仅凭公报中的简短表述，推导每一宗案件的完整证据、具体责任和司法程序结局。台湾中央社报道着重列出遭处分人员及其原任职务；新华社则将会议置于党内纪律、反腐败和治理能力的叙事之中。两种报道选择不同，前者强调权力结构的变化，后者强调组织建设的政策连续性，读者应意识到其材料选择与报道角度。\cite{XinhuaPlenum,CNAPlenum}

\begin{center}\begin{tikzpicture}[>=Stealth,font=\scriptsize,scale=.78,transform shape]
\node[draw=newsblue,rounded corners,fill=blue!5,text width=4.4cm,align=center] (cc) at (0,0) {中央委员会全会\\2026年10月26—29日};
\node[draw=newsblue,rounded corners,fill=blue!4,text width=2.25cm,align=center] (pol) at (-2,-1.45) {政治局\\增选张升民};
\node[draw=newsblue,rounded corners,fill=blue!4,text width=2.25cm,align=center] (rule) at (2,-1.45) {纪律治理决议\\审议通过};
\draw[->,thick] (cc) -- (pol);\draw[->,thick] (cc) -- (rule);
\end{tikzpicture}\par{\footnotesize\color{newsgray}图1｜编辑部原创制度示意；两个结果属不同组织事项，不表示两者互相隶属。依据新华社10月29日全会公报。}\end{center}\color{black}

\newssub{为什么政治局的军方代表值得关注？}政治局参与党的最高层日常决策，中央军委则负责领导军队。高级军官进入政治局，意味着军队系统在党内决策层获得一个明确的正式席位。美联社指出，在此前多名高级军官被撤换后，政治局一度缺乏军方成员；路透社则把张升民此次晋升视为其政治地位的进一步提高。需要避免的误解是：\textbf{一个席位的恢复不等于军队指挥链已经恢复到此前状态}。军委成员数量、岗位分工、军种指挥体系与政治局成员构成属于不同层次，不能由单项任命代替完整的组织结构分析。\cite{APPlenum,ReutersPlenum}

\newssub{「从严治党」意味着什么？}官方公报把纪律建设、监督和反腐败作为长期治理任务，并称要持续加强自我革命。这既是政策宣示，也是对党内干部行为的约束框架。支持者可能认为，持续反腐能减少寻租、强化政策执行与组织纪律；批评者则会追问，集中化的纪律监督如何与透明程序、责任追究和组织内部的信息反馈相协调。两种问题并不互相排斥：即使某次处分有充分事实依据，也不能仅凭处分数量证明制度已达到最佳治理状态；反过来，也不能因为政治斗争的可能性存在，就认定每一起腐败指控必然是虚构的。\cite{XinhuaPlenum,APPlenum}

\newssub{从出席人数能推出什么？}美联社援引分析人士对全会出席比例的计算，认为较低的出席率反映近期人事调整规模较大。但公报提供的是实际出席人数，任何比例都依赖所采用的应到人数、缺席原因和候补递补口径。若不先说明分母，单独拿一个百分比去证明政治清洗的广度，容易把统计结果当成直接证据。更可靠的长期观察方法，是追踪正式名单、机构空缺的填补、职务交接、纪律文件的实施，以及公开政策如何改变军队和地方的决策行为。\cite{XinhuaPlenum,APPlenum}

这次全会的核心新闻并不只是某位将领「升了一级」，而是\textbf{纪律治理成为正式决议、人事调整在党内最高层得到确认，以及军队代表重新进入政治局}三件事同时发生。它们的长期效果尚待后续组织安排与政策执行检验。
}

\newslead{一只「宇宙猫」为何有三张面孔：哈勃与韦布用不同波长观看恒星诞生}{航天科学深度｜NASA／ESA／STScI 10月29日公布；原始观测为2026年5月27日}{

10月29日，NASA与合作机构公开了一组被昵称为「猫头星云」的天文图像。画面中央一条细长、弯曲的暗色尘埃柱，形状近似侧身坐着的猫，亮斑仿佛它的眼睛。这个区域并非新发现的星座，也不是太阳系中的天体，而是距离地球约7500至8500光年的船底座星云的一部分。它真正值得报道的理由，不是恰好赶上万圣节的外形，而是\textbf{哈勃可见光、韦布近红外和中红外成像揭示了同一片物质的不同物理侧面}。\cite{NASACat,NASANIR,NASAMIRI}

\newssub{图像颜色不是肉眼看到的宇宙原色。}望远镜通过不同波段的滤光片记录光子，研究人员把不同波段的单色数据赋予可区分的颜色，再组合成公众看到的彩图。NASA的资料列出，韦布近红外相机NIRCam采用F090W、F187N、F212N、F277W、F335M和F470N等滤光片；中红外仪器MIRI则使用F770W、F1130W和F1500W。不同滤光片的响应范围和强度差异决定了哪些恒星、尘埃与受激发分子更突出。\textbf{假彩色并非伪造数据}，但图像中的蓝、橙、红不能直接等同于天体在肉眼下的颜色或温度。\cite{NASANIR,NASAMIRI}

\begin{center}\begin{tikzpicture}[font=\scriptsize,scale=.8,transform shape]
\foreach \x/\c/\t/\s in {0/blue!25/可见光/尘埃遮挡,2.3/orange!30/近红外/透视部分尘埃,4.6/red!30/中红外/暖尘辐射}{
\fill[\c] (\x,0) rectangle ++(2,1.0);
\draw[gray] (\x,0) rectangle ++(2,1.0);
\node[align=center,text width=1.9cm] at (\x+1,.50) {\t\\\s};
}
\draw[->,thick] (0,-.3)--(6.6,-.3);
\node[below] at (3.3,-.3) {波长由短到长（仅定性示意）};
\end{tikzpicture}\par{\footnotesize\color{newsgray}图2｜原创波段对照，不是NASA观测图。依据NASA／STScI 10月29日公开的Hubble、NIRCam、MIRI说明；实际颜色为滤光片数据映射。}\end{center}\color{black}

\newssub{为什么红外线能「看穿」尘埃？}星际尘埃对短波长光的散射与吸收通常更强，因此哈勃可见光图中的尘埃柱显得较暗，像一块挡住背景的幕布。较长波长的近红外光更容易穿过部分尘埃，韦布因而能显露更多被遮蔽的恒星和复杂表面。中红外则对较暖尘埃和受辐射激发的含碳分子结构敏感；同一条尘埃柱在MIRI图像中又呈现出新的纹理。这里的「看穿」不是无限穿透：足够致密的尘埃在红外也可能不透明，红外亮度还受到发射、吸收与背景光强共同影响。\cite{NASACat,NASANIR,NASAMIRI}

\newssub{猫眼里真有一颗正在诞生的恒星吗？}NASA指出，尘埃柱的明亮「眼睛」可能与一颗包裹在尘埃中的年轻恒星有关，部分明亮边缘可能由附近大质量恒星的紫外辐射照亮、加热并逐步侵蚀。这类过程叫作光致蒸发。恒星诞生不是在静止的云团里悄悄完成的：重力促使局部气体塌缩，附近大质量恒星的辐射与恒星风又能压缩或吹散气体。是否在某一位置确有新生恒星，还需要结合光谱、温度、运动学及多波段证据。官方资料使用「可能」一词，本刊也保留这个证据边界。\cite{NASACat}

\newssub{一张图有三个时间。}韦布所用观测数据采集于5月27日，公众图像发布于10月29日；光从距离数千光年的星云传播到地球还经历了数千年。因此，不能把「10月29日拍到恒星正在诞生」理解成当地那一天发生了新事件。图像处理并不改变光子到达探测器的记录，却可能选择对比度、颜色映射和展示范围。严谨的图注应该同时交代\textbf{观测时间、发布日期、望远镜与波段、处理者和来源}。\cite{NASANIR,NASAMIRI}

这组公开图像提醒我们，天文摄影既是美学作品，也是有校准和选择过程的测量结果。用不同波长研究同一个目标，类似从不同数学投影恢复对象的结构：每种观测丢失一些信息，也揭示另一部分信息。要把漂亮的照片变成物理结论，仍需要光谱模型、误差分析和与其他观测的交叉验证。NASA已公开原始图像页面及完整署名；本刊因无法可靠取得可随源码归档的原始图像文件，改用有明确标注的原创波段示意图，而不以合成画面冒充天文观测。
}

\newsitem{欧洲央行暂停加息：能源价格冲击为何让货币政策左右为难？}{经济政策重点｜欧洲央行10月29日正式决定；路透社市场报道交叉核对}{

欧洲中央银行管理委员会10月29日决定，三项关键政策利率保持不变：存款便利利率2.50\%、主要再融资操作利率2.65\%、边际贷款便利利率2.90\%。这些是欧元体系不同资金操作的利率，不能混写为一个笼统的「欧洲利率」。央行明确表示，能源价格上涨同时提高通胀风险和压低增长前景，政策将逐次会议依据数据决定，不预先承诺未来利率路径。\cite{ECBDecision,ReutersECB}

\begin{center}\begin{tikzpicture}[font=\scriptsize,x=.95cm,y=.44cm]
\foreach \y/\name/\v in {3/存款便利/2.50,2/主要再融资/2.65,1/边际贷款/2.90}{
\node[anchor=east] at (0,\y) {\name};
\fill[newsblue!72] (.15,\y-.24) rectangle ({.15+\v},\y+.24);
\node[anchor=west] at ({.27+\v},\y) {\v\%};
}
\end{tikzpicture}\par{\footnotesize\color{newsgray}图3｜2026年10月29日欧洲央行三项关键利率，单位：年利率百分比；非同比增幅。来源：ECB当日决议；本刊据原始数字绘制。}\end{center}\color{black}

货币政策通常通过调整融资成本影响需求；但能源供应紧张造成的价格上涨首先是一种供给冲击。若央行为抑制通胀而提高利率，居民和企业融资会变贵，需求可能减弱，却不能立即增加石油或天然气供给。若为了增长而降低利率，较宽松的融资环境又可能使能源成本向其他商品与工资传递，形成所谓第二轮通胀效应。央行此次指出，目前工资尚未对能源冲击出现明显反应，基础通胀总体稳定；但冲击的持续时间与规模仍不确定。\cite{ECBDecision}

另一个容易忽略的事实是，政策利率不等于每个国家的实际融资条件。欧元区不同成员国的政府债券利差、银行贷款风险与企业资产负债表不同，同样的政策利率可能产生不同的融资压力。欧洲央行声明还提及其资产购买计划和疫情紧急购买计划的持仓正逐步缩减，意味着货币条件不仅由当天公布的三个利率数字决定。\cite{ECBDecision}

市场据此推测12月可能再次调整利率，但市场价格只是交易者在风险和预期下的综合判断，不能写成央行已决定加息。对未来政策最有解释力的数据，是能源价格持续性、核心通胀、工资增长、信贷与经济活动，而不是单看当天欧元涨跌。
}

\newsitem{朝鲜一次发射约十枚弹道导弹：数量、射程和边境摩擦分别说明什么？}{东北亚安全重点｜韩国军方10月29日通报；对照日本时事社、路透社与美联社}{

韩国联合参谋本部10月29日称，朝鲜从元山附近向东部海域发射约十枚短程弹道导弹，飞行距离约250公里。日本政府表示，初步判断导弹没有落入日本专属经济区，暂未确认附近飞机或船舶受到损害。韩日方面对发射数量与落点的描述基本一致，但这些数字仍属于相关军方初步观测，不等于已公开每枚导弹的型号、弹道高度、载荷或最终落点坐标。\cite{ReutersMissile,JijiMissile,APMissile}

这次试射发生在朝韩边境关系紧张之际。9月非军事区附近地雷爆炸造成三名韩国军人受伤，韩国指责朝方布设，朝鲜否认责任。双方还围绕边境广播、警告射击、铁丝网与军事设施部署发生争执。对这类事件，\textbf{发生了爆炸、有人受伤}与\textbf{究竟是谁布雷}是不同层次的判断：前者可由医疗与现场记录核对，后者需要爆炸物残片、位置、布设时间和独立调查。报道必须将归责留在各自政府的主张层次。\cite{ReutersMissile,APMissile}

从军事技术上看，250公里的飞行距离只能说明这批发射的观测航程，不能直接推导导弹最大射程，也不能证明它们搭载核弹头。弹道导弹的威胁评估还取决于机动性、发射准备时间、命中精度、拦截难度与指挥控制。日本抗议称发射违反联合国安理会相关决议；这是日本政府的法律与安全立场。朝鲜则长期将武器试验解释为防卫和威慑需求。记录各方主张不意味着应把所有指控视为同等程度的已证实事实。\cite{JijiMissile,ReutersMissile}

更值得关注的是冲突管理机制是否还能工作。韩国已提出防止误判的军事接触建议，但朝鲜尚未积极回应。双方在高戒备状态下频繁行动，单次导弹试验未必预示立即开战，却会提高误判、误射和政治升级的风险。
}

\newsitem{数学拾趣｜四色边的完全图：一个Ramsey数上界从62缩到61}{组合数学重点｜Milos Tatarevic，arXiv:2610.36648，10月27日提交；未经同行评议}{

考虑把完全图的每条边涂成红、蓝、绿、黄四种颜色之一。所谓完全图$K_n$，就是$n$个顶点两两相连；所谓\textbf{单色三角形}，就是三个顶点之间的三条边恰好同色。四色Ramsey数$R(3,3,3,3)$是最小的整数$n$，使得不管怎样给$K_n$的边涂这四种颜色，总会出现一个单色三角形。这是一个典型的「无论怎样安排，都躲不开某种结构」的问题。\cite{Tatarevic}

\begin{center}\begin{tikzpicture}[scale=.8, every node/.style={circle,fill=newsblue,inner sep=1.8pt}]
\node (a) at (0,1.25) {}; \node (b) at (-1,0) {}; \node (c) at (1,0) {};
\draw[red,very thick] (a)--(b)--(c)--cycle;
\node[draw=none,fill=none,rectangle,anchor=west,font=\scriptsize] at (1.45,.55) {三条边同色：单色三角形};
\end{tikzpicture}\par{\footnotesize\color{newsgray}图4｜原创定义示意；红色三角形是Ramsey问题要强制出现的结构，不是61顶点证明的实际计算图。}\end{center}\color{black}

一个最简单的例子是只有一种颜色时，$K_3$显然已经包含单色三角形。两种颜色的经典结论$R(3,3)=6$则稍微有趣：在六个人中，无论两人之间的关系按「认识」与「不认识」分成两类，总能找到三个彼此认识或三个彼此不认识的人。证明只需固定一人，看其连向另外五人的边，鸽巢原理保证至少三条同色，再分析那三个端点之间的边。四色问题仍是同一个定义，但组合空间急剧膨胀，不能直接把两色证明机械重复四遍。

Tatarevic的新预印本主张把已知上界从62改进到61，即证明\textbf{任何61个顶点的四色完全图都不可避免有单色三角形}。这是上界的改进，而不是已经求得Ramsey数的精确值。作者的方法包括全局关联计数、对局部邻域结构分类，以及把剩余可能情形转化成布尔可满足性问题（SAT）。SAT求解器可以检查某组逻辑约束是否存在满足赋值；若无解，便排除相应图结构。\cite{Tatarevic}

这里最需要解释的是「计算机证明」的可信边界。论文摘要称，数学归约在Lean中形式化，并用两个SAT求解器排除所有剩余公式。形式化归约能减少人工推导中的遗漏，两个独立求解器也能提供交叉检查；但如果最终的SAT不可满足结论没有被可信内核逐条复核，证明仍依赖求解器实现、输入生成与结果记录。单凭「用了Lean」四个字不能断言所有计算都已由Lean内核验证。本刊未独立运行作者的计算与代码，因此仅报道\textbf{作者提出的61上界证明主张}，不将其作为已完成外部复核的定理。\cite{Tatarevic}

这一研究的魅力在于，61和62只差一个整数，背后却可能需要复杂的组合分类与计算验证。有限数学中的一个单位，往往不是小修小补，而是证明方法能否排除整个庞大构造空间的分界。
}

\newsitem{真核细胞的近亲被培养出来：它们内部简单，却依赖微生物伙伴}{生命科学重点｜《自然》10月28日同行评议论文；Imachi等，DOI:10.1038/s41586-026-11170-7}{

真核细胞拥有细胞核、线粒体等复杂内部结构，动物、植物与真菌都属于真核生物。真核细胞的起源长期与古菌和细菌之间的共生关系有关，但最接近真核生物祖先的古菌究竟长什么样、靠什么生活，仍存在争论。日本等地研究者在《自然》10月28日发表的同行评议论文报告，成功培养出两株属于Heimdallarchaeia类群的古菌HC1与SC1，为这个问题提供了少见的活体实验材料。\cite{NatureArchaea}

两株古菌在显微镜下并没有稳定的真核式内部膜系统，却能形成突出于细胞表面的膜结构与囊泡。它们依赖无氧环境，通过分解多肽获取能量，并与消耗氢气的产甲烷古菌形成互利代谢关系。这里的「共营」意味着，一方代谢产生的氢若不断积累会使反应难以继续，另一方消耗氢后，前者才更容易维持能量代谢。研究还发现，虽然HC1具有与氧利用有关的某些基因，它并不能因此被简单认定为能在有氧条件下独立呼吸。\cite{NatureArchaea,NatureBrief}

论文的推断是，真核生物的古菌近祖先可能内部结构相对简单，具有细胞突出物，并在低氧环境中与细菌祖先形成紧密联系。这支持某些共生起源模型，但\textbf{今天培养出的古菌不是数十亿年前祖先的原样复制品}。从现代亲缘关系反推祖先特征，仍需系统发育树、基因比较、更多培养样本及实验验证。此研究已通过期刊同行评议；这与本期数学预印本尚待外部验证的状态不同。
}

\newsitem{法国零售业遭遇抗议冲击：商店损失与消费信心不能混为一谈}{社会与经济简讯｜路透社10月29日采访法国商会及零售企业}{

法国持续数周的生活成本抗议正在影响年末零售旺季。路透社引述法国商会称，约2400家店铺受损，直接损失超过10亿欧元；若把停业与销售损失计算在内，估计总影响更大。部分商店仍关闭，零售商担忧10月31日及11月5日的新示威进一步影响客流。\cite{ReutersFrance}

必须区分三种数字：被破坏的财产损失、由于营业中断而未实现的销售，以及消费者信心指数。它们不等于同一笔已经确认的净损失，后两项尤其需要模型与调查。抗议者关注物价和生活水平，企业担心就业与经营，政府还要维护公共安全；报道这些利益冲突，不等于为暴力破坏或对示威者的笼统污名化背书。
}

\newsitem{AI代理的「记忆」也要有适用边界：一篇新预印本的启发}{人工智能与软件工程简讯｜Hao等，arXiv:2610.36652，10月27日提交}{

多个AI代理协作时，保存先前任务的经验不一定带来稳定改进：错误规则可能被反复使用，正确规则也可能被用在不适合的场景。Hao等人的DiscernMem预印本提出两种机制：一是比较「使用某条记忆」与「不使用」的执行效果，修订有缺陷的经验；二是记录记忆在不同任务上下文中的收益，学习何时该用、何时不该用。作者在四组基准上报告了性能改进。\cite{DiscernMem}

这是一项作者自报实验结果，并非证明所有代理系统都能通过类似方法获得收益。尤其在软件开发中，旧项目的修复经验若不检查版本、平台和约束条件，可能制造新错误。记忆系统的关键不只是储存容量，更是可检验的适用条件与撤销机制。
}



\newsitem{红外相机能直接「长」在芯片上吗？《自然》报道室温水溶液生长硫化铅探测器}{材料、化学与电子工程重点｜李森、刘静、高亮等；《自然》10月28日发表；DOI:10.1038/s41586-026-11222-y}{
红外成像设备需要把不可见的红外光转化为电信号。传统制造方案通常先制作探测器阵列，再与负责读取信号的集成电路连接；这种混合集成方法成熟，但在材料对准、连接密度、成本和制造良率方面存在工程挑战。华中科技大学等机构的研究者在《自然》10月28日发表论文，报告一种\textbf{在室温水溶液中直接于读出芯片上生长硫化铅（PbS）探测层}的方法，称为单片原位集成（MISI）。这是经过同行评议的实验研究，但其工程化成本与长期可靠性仍需进一步检验。\cite{NaturePbS}

硫化铅是一种对红外光有响应的半导体。论文的关键难点并不只是「把PbS沉积上去」，而是让生长过程与已有芯片兼容，同时得到电子质量足够高、各像素足够一致的薄膜。研究团队利用铅与氢氧根离子的配位平衡，控制水溶液中的生长活性物种，并尽量抑制溶液中自行形成颗粒的均相成核，让材料优先在目标表面结晶。研究报告载流子迁移率为124\,cm$^2$\,V$^{-1}$\,s$^{-1}$，载流子浓度约$2\times10^{15}$\,cm$^{-3}$；这些是实验样品的材料指标，不能直接当成整机的图像质量分数。\cite{NaturePbS}

\begin{center}\begin{tikzpicture}[>=Stealth,font=\scriptsize,scale=.78,transform shape]
\node[draw=newsblue,rounded corners,fill=blue!5,minimum width=1.7cm,minimum height=.75cm,align=center] (r) at (0,0) {读出电路\\ROIC};
\node[draw=newsblue,rounded corners,fill=blue!5,minimum width=1.8cm,minimum height=.75cm,align=center] (s) at (3,0) {室温溶液\\原位生长PbS};
\node[draw=newsblue,rounded corners,fill=blue!5,minimum width=1.6cm,minimum height=.75cm,align=center] (d) at (6,0) {红外像素\\直接集成};
\draw[->,thick] (r)--(s);\draw[->,thick] (s)--(d);
\end{tikzpicture}\par{\footnotesize\color{newsgray}图5｜本刊原创工艺流程示意，非论文原始显微照片。来源：Li等，《自然》，2026年10月28日，DOI:10.1038/s41586-026-11222-y。}\end{center}\color{black}

通过锑掺杂，研究者在2.5微米波长报告比探测率$4.4\times10^{11}$ Jones，并制作出640$\times$512像素的红外焦平面阵列。比探测率综合反映探测器在噪声背景下识别弱光信号的能力，但其数值与测量波长、温度、带宽和测试条件有关。论文声称在该特定条件下超过商用InGaAs探测器，并展示了清晰的红外成像；\textbf{这不意味着在所有波段、环境和使用寿命下都全面优于商业器件}。\cite{NaturePbS}

研究还把类似工艺应用于柔性读出电路，并展示了向硒化铅材料扩展的可能。材料科学的进步往往不只体现在单个像素的灵敏度，而在于能否把可控结晶、芯片兼容、均匀性和批量制造连成完整工艺链。未来仍需观察长期热稳定性、器件封装、环境耐受性、批次良率，以及铅化合物生产和废弃处理的环境安全。本文不把「室温制备」误写为「整个产品零能耗制造」；它只是降低了特定沉积步骤的温度需求。
}

\newsitem{莫斯科举办俄非峰会：贸易、粮食与安全合作背后，是非洲国家的选择权}{国际外交重点｜10月28—29日俄罗斯—非洲峰会；美联社连续报道}{
俄罗斯总统普京在莫斯科举行的俄非峰会上呼吁深化政治、经济和安全合作。美联社报道，约48个非洲国家派出代表，其中约20个国家的领导人出席。普京将俄罗斯与非洲的关系放在「多极世界」与反殖民历史的叙事中，支持提高非洲在联合国安理会等机构中的代表性，并称俄罗斯对非贸易、粮食与化肥出口正在扩大。这些贸易和粮食数字主要来自普京本人发言，\textbf{不能不加说明地当作独立审计后的双边统计}。\cite{APRussiaAfrica1,APRussiaAfrica2}

峰会的政治意义不能只从俄罗斯的角度理解。非洲有54个主权国家，资源禀赋、殖民历史、政体和外交利益差异巨大。对部分国家来说，与俄罗斯开展能源、矿产或教育合作，是在欧美、中国、海湾国家及其他伙伴之间增加选择；对另一些国家而言，俄罗斯的军售与安全承包服务可能与本国政权安全直接相连。把非洲各国一概称为「俄罗斯阵营」，既忽略国家之间的差异，也抹去了这些国家自主进行外交权衡的能力。\cite{APRussiaAfrica1}

安全合作尤其需要谨慎核查。美联社指出，俄罗斯在萨赫勒地区的军事影响力与原瓦格纳集团及后来的「非洲军团」相关。支持者认为这些力量可帮助政府打击武装组织，批评者则关注平民伤亡、军事政府的权力巩固及问责缺失。军事合作协议的数量、实际部署人员、合同条款和战场结果不是同一指标；一国签署合作协议，也不意味着它的民众普遍支持某种外交路线。\cite{APRussiaAfrica2}

俄罗斯希望以粮食、化肥、矿产投资和外交支持拓展伙伴关系，但峰会声明距离可验证的项目落地仍有距离。后续应追踪具体合同、贷款条件、港口和运输项目、就业与当地财政收入，尤其观察利益是否流向普通居民。国际政治的「影响力」不是抽象的国旗数量，而是持续合作是否能经受经济与安全风险的考验。
}

\newsitem{英国推出「延续性预算」：税收、能源补贴和财政余量如何同时计算？}{欧洲经济简讯｜路透社10月29日市场报道；政策细则仍需逐项核查}{
英国财政大臣John Healey公布新一轮预算，涉及银行、博彩、储蓄、房东等领域税负调整，以及在线市场增值税规则和居民能源账单缓冲措施。路透社报道，政府用于满足主要财政规则的预测余量从240亿英镑降至150亿英镑；英镑当天变化不大。\cite{ReutersUKBudget}

「财政余量」是预测收入、支出和财政规则之间的空间，不是国库里现成的150亿英镑现金；而汇率没有剧烈波动，也不能证明预算对所有家庭和行业都没有影响。需要等待完整预算文件和分配分析，才能判断不同收入群体、租户、企业与金融机构实际承担的成本。英国能源政策还受到中东战争引发的能源价格波动影响，短期补贴与长期财政可持续性必须分别衡量。
}

\par\medskip{\large\bfseries\color{newsblue}参考资料与编辑说明}\par
{\scriptsize 本期时政每日必检；数学核查arXiv数学新列表及相关预印本。政治人物、机构及军方的声明均按来源归属。科学原始图片可在NASA页面核查；本期因无法取得可归档的原图文件，采用三张原创TikZ图和一张原创数据图，不将示意图伪装为真实观测。图表的数字、口径及来源见图注。}\par
\begingroup\fontsize{6}{7.1}\selectfont
\renewcommand{\refname}{参考文献}
\bibliographystyle{plainurl}\bibliography{refs}
\endgroup
\end{documents}
\end{document}
